\textbf{1. Kiến thức cần nhớ}
\begin{itemize}
\item \textbf{Phương trình bậc nhất hai ẩn} $x$ và $y$ là hệ thức dạng:
$$\boxed{ax + by = c}$$
trong đó $a, b, c$ là các số đã biết ($a$ và $b$ không đồng thời bằng $0$).
\item \textbf{Nghiệm của phương trình bậc nhất hai ẩn}: Cặp số $(x_0; y_0)$ thỏa mãn $ax_0 + by_0 = c$ được gọi là một nghiệm của phương trình.
\begin{itemize}
\item Phương trình bậc nhất hai ẩn luôn có \textbf{vô số nghiệm}.
\item Tập nghiệm của phương trình $ax + by = c$ được biểu diễn bởi đường thẳng $d\colon ax + by = c$ trên mặt phẳng tọa độ $Oxy$.
\end{itemize}
\item \textbf{Hệ hai phương trình bậc nhất hai ẩn} có dạng:
$$\boxed{\begin{cases} ax + by = c \\ a'x + b'y = c' \end{cases}}$$
(với $a^2 + b^2 \ne 0$ và $a'^2 + b'^2 \ne 0$).
\item \textbf{Số nghiệm của hệ}: Số nghiệm của hệ bằng số giao điểm của hai đường thẳng $d\colon ax + by = c$ và $d'\colon a'x + b'y = c'$:
\begin{itemize}
\item Hệ có nghiệm duy nhất $\iff d$ cắt $d'$ ($\dfrac{a}{a'} \ne \dfrac{b}{b'}$ khi các mẫu khác $0$).
\item Hệ vô nghiệm $\iff d \parallel d'$ ($\dfrac{a}{a'} = \dfrac{b}{b'} \ne \dfrac{c}{c'}$).
\item Hệ có vô số nghiệm $\iff d \equiv d'$ ($\dfrac{a}{a'} = \dfrac{b}{b'} = \dfrac{c}{c'}$).
\end{itemize}
\end{itemize}

\textbf{2. Phương pháp giải các dạng}

%% DANG 9D1?1-1
\textbf{Dạng 1. Nhận biết phương trình bậc nhất hai ẩn, nghiệm của phương trình}
\begin{enumerate}
\item Kiểm tra dạng $ax + by = c$ với $a, b$ không đồng thời bằng $0$ và các biến $x, y$ có số mũ là $1$.
\item Để kiểm tra xem cặp số $(x_0; y_0)$ có là nghiệm không, ta thay $x = x_0$ và $y = y_0$ vào vế trái: nếu $VT = VP = c$ thì là nghiệm; ngược lại thì không phải nghiệm.
\item Để tìm nghiệm tổng quát: biểu diễn ẩn này theo ẩn kia, ví dụ: $y = \dfrac{c - ax}{b}$ (với $b \ne 0$) hoặc $x = \dfrac{c - by}{a}$ (với $a \ne 0$).
\end{enumerate}

%% DANG 9D1?1-2
\textbf{Dạng 2. Biểu diễn tập nghiệm trên mặt phẳng tọa độ}
\begin{enumerate}
\item Tìm hai điểm phân biệt thuộc đường thẳng $d\colon ax + by = c$:
\begin{itemize}
\item Cho $x = 0 \implies y = \dfrac{c}{b}$ (nếu $b \ne 0$).
\item Cho $y = 0 \implies x = \dfrac{c}{a}$ (nếu $a \ne 0$).
\end{itemize}
\item Xác định hai điểm đó trên mặt phẳng tọa độ $Oxy$ rồi vẽ đường thẳng đi qua hai điểm.
\end{enumerate}

%% DANG 9D1?1-3
\textbf{Dạng 3. Nhận biết hệ hai phương trình bậc nhất hai ẩn, kiểm tra nghiệm của hệ}
\begin{enumerate}
\item Kiểm tra xem hệ có đúng gồm hai phương trình bậc nhất hai ẩn cùng chứa hai biến $x, y$ hay không.
\item Để kiểm tra cặp số $(x_0; y_0)$ có là nghiệm của hệ không: thay $(x_0; y_0)$ vào \textbf{cả hai phương trình}. Cặp số là nghiệm của hệ khi và chỉ khi nó là nghiệm của đồng thời cả hai phương trình.
\end{enumerate}
